%------------------------------------------------------------------------------%
\section{Inequações Exponenciais}

Para resolvermos uma Inequação exponencial do tipo
\[
  a^{f(x)} \leq a^{g(x)} \qquad\text{ ou }\qquad
  a^{f(x)} \geq a^{g(x)},
\]
devemos lembrar que a função $y=a^x$ é crescente se $a>1$ e decrescente
se $0 < a< 1$.
Portanto, antes de tudo, observamos a situação das bases dos dois membros e,
caso as bases sejam diferentes, reduza-as as mesmas bases e depois monte
uma inequação respeitando as seguintes condições:
\begin{enumerate}
    \item Se $a>1$ conserve o sinal original da inequação. Isso se dá pelo fato
      da função exponencial ser crescente se $a>1$.

    \item Se $0< a< 1$ inverta o sentido da desigualdade. Essa fato, se dá por
      conta da função exponencial ser decrescente quando $0< a< 1$.
\end{enumerate}

%------------------------------------------------------------------------------%
\begin{example}
  Resolva a inequação exponencial
  \(
    \left( \frac{1}{2} \right)^{2x+2} <
    \left( \frac{1}{2} \right)^{x^2-1}
  \).

  \solution
  Note que nesse caso temos uma inequação exponencial e como as bases de ambos
  os membros são iguais
  \[
    a = \frac{1}{2}
  \]
  basta montarmos uma inequação com os expoentes,
  invertendo o sentido da desigualdade, visto que $0 < a < 1$.

  Assim, precisamos resolver a seguinte inequação exponencial
  \begin{align*}
    2x + 2       & > x^2 - 1 \\
    x^2 - 2x - 3 & < 0,
  \end{align*}
  que é uma inequação do segundo grau, cuja solução é
  \[
    S = \left\{ x \in \R \st -1 < x < 3 \right\}.
  \]
\end{example}
%------------------------------------------------------------------------------%

%------------------------------------------------------------------------------%
\begin{example}
  Resolva a inequação exponencial \(2^{x+2} \geq 64\).

  \solution
  Inicialmente, observamos que $64=2^6$.
  Portanto, estamos interessado em resolver a inequação
  \[
    2^{x+2} \geq 2^6.
  \]
  Como já deixamos ambos os membros da desigualdade na mesma base
  então montamos uma inequação com os expoentes mantendo o sentido
  da desigualdade original, uma vez que a base $a = 2 > 1$.

  Assim, temos que resolver a seguinte inequação do primeiro grau
  \begin{align*}
    x+2 & \geq 6 \\
    x   & >    4.
  \end{align*}
  Logo, a solução da inequação pedida é
  \[
    S = \left\{ x \in \R \st x > 4 \right\}.
  \]
\end{example}
%------------------------------------------------------------------------------%

%------------------------------------------------------------------------------%
\begin{example}
  Determine o domínio $D$ da função \(f(x) = \sqrt{5^x-125}\).

  \solution
  A fim de que exista $f(x)$, devemos ter
  \[
    5^x-125 \geq 0,
  \]
  uma vez que, não existe raiz quadrada de número negativo. Assim,
  \begin{align*}
    5^x-125 & \geq 0   \\
    5^x     & \geq 125 \\
    5^x     & \geq 5^3 \\
    x       & \geq 3
  \end{align*}
  Assim, o domínio da função é
  \[
    D = \{ x \in \R \st x \geq 3 \}.
  \]
\end{example}
%------------------------------------------------------------------------------%
